\subsection*{A stress-cone criterion for the two corrected azimuthal stages}

The Appendix-A condition is not merely a sign restriction on
$\ell=\partial_{\log X}\log(\sqrt X E)$. The source's equations
(4.11), (4.20)--(4.22) give
\[
a=2-2\ell,\quad b_s=2D_XU/E,\quad
p_1=XQ_s/L,\quad w=N_s/(EQ_s).
\]
For $a>0$ put $t_s=-b_s/a$,
$v_s=a+b_s^2/a$, $P_c=p_1(1-b_sw/a)$ and
$J_c=p_1(w+b_s/a)$. When $v_s>2$, the full cone reads
$P_c>v_s$ and $(v_s-2)J_c^2<2(P_c-v_s)^2$.

\begin{proposition}[Uniform corrected-support sufficient condition]
\label{prop:A11-cone-source}
On the two genuine (A.11) azimuthal-bump supports,
$U=M=J=0$ and hence $b_s=0$. Assume the corrected slope
$-3\lambda/2\le\ell\le-\lambda/2$, $0<\lambda\le1/20$.
If throughout those supports
\[
Q_s\ge\lambda/4,\quad |N_s|\le E\lambda/4,
\quad X\ge18/\lambda,
\]
then $a>2$, $p_1\ge2a$ and the normalized cone margin obeys
\[
2(1-a/p_1)^2-(a-2)w^2\ge7/20.
\]
\end{proposition}
\begin{proof}
The slope gives $2+\lambda\le a\le2+3\lambda\le43/20$.
As $L\le1$, $p_1=XQ_s/L\ge9/2>2a$.
Also $|w|\le1$. Therefore
$2(1-a/p_1)^2-(a-2)w^2\ge1/2-3/20=7/20$.
Multiply by $p_1^2>0$ and invoke (4.22).
\end{proof}

\begin{proposition}[Conditional source lower bound for $Q_s$]
\label{prop:A11-cone-Qs}
Retain the previously proved A.11 bounds for $c$, $c_\eta$,
$\ell_c+\lambda$, $\partial_\eta\ell_c$ and the initial
first-bump discrepancy:
\[
\|c\|\le192\lambda^{28},\quad
\|c_\eta\|<12\lambda^{28},\quad
|\ell_c+\lambda|<76800000\lambda^{28},\quad
|\ell_{c,\eta}|<5000000\lambda^{28},
\]
and at the first bump center
$|r-(1-\lambda)^{-1}|<64\lambda^{28}$,
$|r_\eta|<4\lambda^{28}$ for $r=I/(XH)$.
If $0<h\le\lambda/100$, then $Q_s\ge\lambda/2$ on both
supports and the intervening constant-slope interval.
\end{proposition}
\begin{proof}
Moving to the left endpoint of the first support enlarges
the initial discrepancies to $80\lambda^{28}$ and
$5\lambda^{28}$. The total interval is shorter than three
logarithmic units. Solve the damped equation
$r'+(1+\ell_c)r=1$, using $0\le r\le20/17$:
\[
|r-(1-\lambda)^{-1}|<3\cdot10^8\lambda^{28},
\qquad |r_\eta|<2\cdot10^7\lambda^{28}.
\]
On either correction support $E=E_{\rm base}(1+c_i\beta_i)$,
with $E_{\rm base}$ independent of $\eta$.
Since $\|\beta_i\|_\infty\le640/3$,
$|H_\eta/H|\le5120\lambda^{28}$, giving
$|I_\eta/(XH)|<2.1\cdot10^7\lambda^{28}$.
The exact identity (4.16) at $U=M=J=0$ is
\[
Q_s=-1+(1-h)r-D\eta I_\eta/(XH).
\]
Hence $Q_s\ge(\lambda-h)/(1-\lambda)-
3.2\cdot10^8\lambda^{28}\ge\lambda/2$, using
$3.2\cdot10^8(1/20)^{27}<49/100$.
\end{proof}

The remaining sufficient support conditions thus reduce to
$|N_s|\le E\lambda/2$ and $X\ge9/\lambda$.
The radial condition can be imposed by increasing $X_R$
after the other source parameters. The first condition is an
unverified source-specific interval requirement on the actual
corrected $N_s$.

From the source's (4.16), on these supports
\[
Q_s=-1+\frac{(1-h)I-D\eta I_\eta}{XH},\quad
N_s=\frac{4h\eta S-(1-\eta^2)S_\eta}{X}
 +4A\eta\Pi-(1-\eta^2)\Pi_\eta .
\]
Thus a genuine interval certificate needs the six cumulative
moment/pressure fields and their indicated $\eta$ derivatives.

\begin{lemma}[Quantitative stability of the A.24 margins]
If both the uncorrected and corrected triples $(a,b_s,w)$
satisfy $a\ge a_0>0$, $|a|\le A_0$, $|b_s|\le B_0$,
$|w|\le W_0$ and corresponding component differences
at most $\epsilon$, then for
$m_1=a-b_sw$ and
$m_2=2-2b_sw-b_s^2/a-(a-2)w^2$,
\[
|\Delta m_1|\le(1+B_0+W_0)\epsilon,
\]
\[
|\Delta m_2|\le
\left[2(B_0+W_0)+2B_0/a_0+B_0^2/a_0^2+
W_0^2+2(A_0+2)W_0\right]\epsilon.
\]
\end{lemma}
\begin{proof}
Expand the products, use
$|\Delta(b_s^2)|\le2B_0\epsilon$,
$|\Delta(a^{-1})|\le\epsilon/a_0^2$, and
$|\Delta(w^2)|\le2W_0\epsilon$.
\end{proof}

On the earlier joining stages the source requires only the
relaxed cone (4.21); the full condition $v_s>2$ applies on
the designated outer intervals. Strict pulse margins from
Appendix A.5 require separate interval bounds on the
source's $o(1)$ errors. The exact checker supplied with this
insert verifies scalar cone implications from externally
certified input intervals, not the actual exterior schedule.
The complete source stress cone and the proposed unforced
Theorem~1.1 remain unproved.
